\section{Prescribed solutions and their uniqueness}\label{sec:analytic}
A construction below first specifies a smooth velocity $U$ and then defines its residual force. We need uniqueness only for those data. On a torus, the comparison class consists of classical velocities $v$ with $v,v_t$ and spatial derivatives through order two continuous on every finite closed time cylinder, and periodic pressures $p$ with $p,\nabla p$ continuous and mean zero. On $\mathbb R^3$, we use smooth classical solutions satisfying, for each $T<\infty$,
\begin{equation}\label{eq:comparison}
 v\in C([0,T];H^2)\cap C^1([0,T];L^2),\qquad
 p\in C([0,T];H^1),\qquad
 \sup_{[0,T]\times\mathbb R^3}(|v|+|\nabla v|)<\infty.
\end{equation}
These conditions concern competitors as well as the exhibited solution. In the whole-space constructions the reference velocity has fixed compact support; a competitor is not required to have that support.

\begin{lemma}[Residual realization]\label{lem:b-comparison}
Fix $\nu>0$ and let $U$ be a prescribed smooth divergence-free velocity with $U(0)=0$, bounded together with its spatial gradient on every finite time interval. In the whole-space case assume also the velocity conditions in \eqref{eq:comparison}. Then the force
\[
 f=\mathcal F_\nu[U]:=U_t+(U\cdot\nabla)U-\nu\Delta U
\]
has $(U,0)$ as its unique solution in the stated comparison class. Its material trajectories exist uniquely for all finite forward times and give smooth diffeomorphisms at each such time. A mean-zero velocity on the torus has a mean-zero residual force.
\end{lemma}
\begin{proof}
Substitution proves existence. If $(v,p)$ is another solution, put $w=v-U$. The difference equation is
\[
 w_t+(v\cdot\nabla)w+(w\cdot\nabla)U
       =\nu\Delta w-\nabla p,\qquad w(0)=0,\quad\operatorname{div}w=0.
\]
On the torus, multiply by $w$ and integrate periodically. Transport by $v$ and pressure give zero, and diffusion gives $-\nu\|\nabla w\|_2^2$. The stated classical regularity justifies these integrations and time differentiation.

On $\mathbb R^3$, insert a smooth cutoff $\chi_R$ equal to one on the ball of radius $R$, supported in the ball of radius $2R$, and with $|\nabla\chi_R|\le C/R$. The boundary errors in the transport, pressure and diffusion integrations are bounded respectively by
\[
 CR^{-1}\|v\|_\infty\|w\|_2^2,\qquad
 CR^{-1}\|p\|_2\|w\|_2,\qquad
 CR^{-1}\|\nabla w\|_2\|w\|_2.
\]
They tend to zero. The $C^1(L^2)$ condition justifies differentiation of the squared norm; $C(H^2)$ and the finite-interval bounds justify the limiting products and integrations. Equivalently one can pass to the integrated-in-time identity, with the same estimates uniform on each finite interval. In either domain the result is
\[
 \frac12\frac{d}{dt}\|w\|_2^2+\nu\|\nabla w\|_2^2
 \le\|\nabla U\|_{\infty,\mathrm{op}}\|w\|_2^2.
\]
Gronwall's inequality gives $w=0$. The equation then gives $\nabla p=0$. Mean-zero normalization on the torus fixes $p=0$; on $\mathbb R^3$, the only spatially constant $H^1$ pressure is zero.

Bounded speed prevents finite-time escape of a material trajectory. Local spatial Lipschitz continuity gives uniqueness, and solving backwards on a finite interval gives the inverse flow. Smooth dependence on the initial point gives its smoothness. Finally, the torus integral of $\Delta U$ and of $(U\cdot\nabla)U=\operatorname{div}(U\otimes U)$ vanishes, proving the mean assertion.
\end{proof}

The lemma is the classical difference-energy argument; it is not a general large-data existence theorem. All forces below come with their explicitly constructed global smooth solution. A smooth compactly supported velocity with the derivative bounds we establish automatically belongs to the whole-space velocity class in \eqref{eq:comparison}.
